% vim: spl=pt
\begin{exercício}{Propriedades de operadores limitados e normais}{ex3}
  Seja \(T, S \in \bounded(\hilbert)\) operadores limitados. Mostre que
  \begin{enumerate}[label=(\alph*)]
    \item \(\norm{TS} \leq \norm{T}\norm{S}\); e
    \item \(\ker{T^*} = \ran{T}^\perp\).
  \end{enumerate}
  Um operador \(T\) é dito normal se \(T T^* = T^* T.\) Mostre que para um operador normal \(T \in \bounded(\hilbert)\) é verdade que
  \begin{enumerate}[label=(\alph*), start=3]
    \item \(\norm{Tx} = \norm{T^*x},\) para todo \(x \in \hilbert;\)
    \item \(\ker{T^*} = \ker{T}\);
    \item \(\norm{(T^*T)^n} = \norm{(T^n)^*T^n} = \norm{T^n}^2\);
    \item se \(x \in \hilbert\) e \(\lambda \in \mathbb{C}\) são tais que \(Tx = \lambda x,\) então \(T^*x = \conj{\lambda} x\); e
    \item autovetores de \(T\) com autovalores distintos são ortogonais.
  \end{enumerate}
\end{exercício}
\begin{proof}[Resolução de (a)]
  Temos
  \begin{equation*}
    \norm{TS} = \sup_{\norm{x} = 1}{\norm{TSx}} \leq \sup_{\norm{x} = 1}{\norm{T} \norm{Sx}} = \norm{T} \norm{S},
  \end{equation*}
  como desejado.
\end{proof}
\begin{proof}[Resolução de (b)]
  Temos
  \begin{align*}
    y \in \ker T^* &\iff T^*y = 0\\
                   &\iff \forall x \in \hilbert, \inner{T^*y}{x} = 0\\
                   &\iff \forall x \in \hilbert, \inner{y}{Tx} = 0\\
                   &\iff \forall x \in \ran T, \inner{y}{x} = 0\\
                   &\iff y \in \ran T^\perp,
  \end{align*}
  portanto \(\ker T^* = \ran T^\perp.\) Vale notar que utilizamos que \(T\) é limitada, portanto \(T^{**} = T\).
\end{proof}
\begin{proof}[Resolução de (c)]
  Como \(T\) é normal, temos
  \begin{align*}
    \norm{Tx}^2 &= \inner{Tx}{Tx}\\
                            &= \inner{T^*Tx}{x}\\
                            &= \inner{T T^*x}{x}\\
                            &= \inner{T^*x}{T^*x}\\
                            &= \norm{T^*x}^2
  \end{align*}
  para todo \(x \in \hilbert.\) Assim, \(\norm{Tx} = \norm{T^*x}\) para todo \(x \in \hilbert.\)
\end{proof}
\begin{proof}[Resolução de (d)]
  Temos
  \begin{equation*}
    x \in \ker T \iff \norm{Tx} = 0 \iff \norm{T^*x} = 0 \iff x \in \ker{T^*},
  \end{equation*}
  portanto \(\ker{T} = \ker{T^*}.\)
\end{proof}
\begin{proof}[Resolução de (e)]
  Como \(T\) é normal, temos
  \begin{equation*}
    (T^*)^n T^n = (T^*T)^n,
  \end{equation*}
  já que \(T\) e \(T^*\) comutam. Desse modo, como \(\norm{A}^2 = \norm{A^*A}\) para todo operador limitado \(A \in \bounded(\hilbert),\) segue que
  \begin{equation*}
    \norm{T^n}^2 = \norm{(T^n)^*(T^n)} = \norm{(T^*T)^n},
  \end{equation*}
  como desejado. Utilizamos aqui que se \(A \in \bounded(\hilbert),\) então \(A^m \in \bounded(\hilbert)\) para todo \(m \in \mathbb{N},\) já que \(\norm{A^mx} \leq \norm{A}^m \norm{x}\) para todo \(x \in \hilbert.\)
\end{proof}
\begin{proof}[Resolução de (f)]
  De (d) segue que
  \begin{equation*}
    x \in \ker(T - \lambda) \iff x \in \ker(T^* - \conj\lambda),
  \end{equation*}
  portanto, se \(x\) é autovetor de \(T\) associado ao autovalor \(\lambda,\) então \(x\) é autovetor de \(T^*\) associado ao autovalor \(\conj\lambda\).
\end{proof}
\begin{proof}[Resolução de (g)]
  Sejam \(x_1, x_2 \in \hilbert\setminus\set{0}\) autovetores de \(T\) com autovalores \(\lambda_1, \lambda_2 \in \mathbb{C}\) tais que \(\lambda_1 \neq \lambda_2.\) Notemos que, por (f),
  \begin{equation*}
    \lambda_1 \inner{x_1}{x_2} = \inner{T^*x_1}{x_2} = \inner{x_1}{Tx_2} = \lambda_2 \inner{x_1}{x_2},
  \end{equation*}
  portanto
  \begin{equation*}
    (\lambda_1 - \lambda_2) \inner{x_1}{x_2} = 0.
  \end{equation*}
  Desse modo, como \(\lambda_1 - \lambda_2 \neq 0,\) concluímos que
  \begin{equation*}
      i \neq j: \ker(T - \lambda_i) \subset \ker(T - \lambda_j)^\perp,
  \end{equation*}
  isto é, autovetores com autovalores distintos são ortogonais.
\end{proof}
